\documentclass[journal]{IEEEtran}

% Packages
\usepackage{cite}
\usepackage{amsmath,amssymb,amsfonts,amsthm}
\usepackage{algorithm}
\usepackage{algorithmic}
\usepackage{graphicx}
\usepackage{textcomp}
\usepackage{xcolor}
\usepackage{siunitx}
\usepackage{booktabs}
\usepackage{multirow}
\usepackage{array}
\usepackage{float}
\usepackage{subcaption}
\usepackage{url}
\graphicspath{{./figures/}}

% Theorem environments
\newtheorem{theorem}{Theorem}
\newtheorem{corollary}[theorem]{Corollary}
\newtheorem{proposition}[theorem]{Proposition}
\newtheorem{lemma}[theorem]{Lemma}
\newtheorem{definition}{Definition}

\def\BibTeX{{\rm B\kern-.05em{\sc i\kern-.025em b}\kern-.08em
    T\kern-.1667em\lower.7ex\hbox{E}\kern-.125emX}}

\begin{document}

\title{Quantum-Certified Anonymization: Irreversibility Beyond Computational Hardness}

\author{Daniel~Mo~Houshmand%
\thanks{D. M. Houshmand is with QDaria AS, Oslo, Norway (e-mail: mo@qdaria.com). ORCID: 0009-0008-2270-5454.}%
\thanks{Patent pending: Norwegian Industrial Property Office (Patentstyret), filed March 2026.}}

\markboth{Proceedings on Privacy Enhancing Technologies, 2026}%
{Houshmand: Quantum-Certified Anonymization}

\maketitle

\begin{abstract}
We present the first data anonymization system whose irreversibility is guaranteed by the Born rule of quantum mechanics rather than by computational hardness assumptions. Existing anonymization techniques, including $k$-anonymity, $\ell$-diversity, differential privacy, and PRNG-based one-time pad masking, derive their security from the computational infeasibility of recovering the random seed used during the anonymization transformation. An adversary who captures the PRNG state through memory forensics, side-channel attacks, or insider access can reconstruct every ``random'' value and reverse the anonymization completely. We introduce \textsc{QRNG-OTP-Destroy}, a protocol that replaces each personally identifiable value with a token derived from quantum random numbers produced by measuring qubits in superposition, then irreversibly destroys the mapping. Because quantum measurement outcomes are governed by the Born rule and no deterministic seed exists, the anonymization is information-theoretically irreversible: no adversary, regardless of computational power (classical or quantum), can recover the original data. We formalize three tiers of irreversibility (computational, information-theoretic, and physics-guaranteed), prove that no classical PRNG-based method can achieve the strongest tier, and prove that \textsc{QRNG-OTP-Destroy} does. We report on an implementation with 10 progressive anonymization levels, backed by multi-provider quantum entropy from IBM Quantum (156-qubit processors), Rigetti, and qBraid, validated by 109 unit and integration tests. The system is the first to provide an auditable chain from quantum hardware provenance to GDPR Recital~26 compliance.
\end{abstract}

\begin{IEEEkeywords}
Anonymization, quantum random number generation, Born rule, differential privacy, GDPR, information-theoretic security, one-time pad
\end{IEEEkeywords}

%% ====================================================================
\section{Introduction}
\label{sec:intro}
%% ====================================================================

The dominant threat model in data privacy has shifted. Intelligence agencies and well-resourced adversaries now routinely execute harvest-now, decrypt-later (HNDL) strategies: collecting encrypted and anonymized datasets today with the expectation that advances in computing will render current protections reversible in the future. For encrypted data, the response has been migration to post-quantum cryptographic algorithms standardized by NIST in August 2024 (FIPS~203, FIPS~204, FIPS~205). For anonymized data, no equivalent migration path exists. The anonymization community has not reckoned with the fact that every deployed anonymization tool derives its irreversibility from the same computational hardness assumptions that the cryptographic community has already judged insufficient.

Consider the operational reality. Every anonymization system in production today, from academic tools such as ARX~\cite{prasser2014arx} and sdcMicro~\cite{templ2017sdc} to industrial systems including Google's differential privacy library~\cite{wilson2020dpsql}, Apple's local differential privacy~\cite{apple2017dp}, and the OpenDP framework, uses a classical pseudo-random number generator (PRNG) or cryptographically secure PRNG (CSPRNG) as its entropy source. A CSPRNG, whether ChaCha20, AES-CTR-DRBG, or the Linux kernel's \texttt{/dev/urandom}, is deterministic: given the internal state, every output can be reproduced. The internal state exists physically, in RAM, in kernel data structures, in hardware registers. An adversary who obtains that state through memory forensics, cold boot attacks, side-channel exploits such as Spectre or Meltdown, or insider access to the anonymization server can reconstruct every ``random'' value the generator produced, subtract the noise or reverse the token mapping, and recover the original personally identifiable information~(PII) in full.

This is not a theoretical concern. Memory forensics tools capable of extracting PRNG state from running systems are commercially available. Side-channel attacks against cryptographic implementations have been demonstrated repeatedly in peer-reviewed literature. The HNDL adversary does not need to break the PRNG algorithm; they need only capture its state at the moment of anonymization. Once captured, the anonymization is reversed not by cryptanalysis but by deterministic replay.

The European Union's General Data Protection Regulation (GDPR) draws a sharp legal distinction in Recital~26 between anonymous data and pseudonymous data. Anonymous data, information that ``does not relate to an identified or identifiable natural person,'' falls entirely outside the regulation's scope. Pseudonymous data, where re-identification remains possible using additional information, remains personal data subject to the full weight of data protection obligations. The economic and operational difference between these two categories is substantial. Organizations that can demonstrate true anonymization can share, archive, and process data without consent requirements, data subject access requests, or retention limits.

The gap in the current framework is this: no existing anonymization method can prove that re-identification is impossible as a matter of physical law. Every method proves, at best, that re-identification is computationally infeasible under stated assumptions. Those assumptions may hold today. They may not hold in a decade. They do not hold against an adversary who captures the PRNG state.

This paper presents the first anonymization system where irreversibility is guaranteed by the Born rule of quantum mechanics. The system, called \textsc{QRNG-OTP-Destroy}, replaces each PII value with a token derived from quantum random numbers generated by measuring qubits in superposition on quantum computing hardware, then securely destroys the mapping between original values and replacement tokens. Because the quantum measurement outcomes that produced the tokens are governed by the Born rule, there is no seed, no hidden state, and no deterministic reconstruction path. The mapping's destruction eliminates the only artifact that could link tokens to original values. The result is anonymization whose irreversibility does not depend on any computational hardness assumption and holds regardless of advances in classical computing, quantum computing, or the resolution of the P versus NP problem.

Our contributions are:

\begin{enumerate}
\item \textbf{Formal definitions.} We define three tiers of anonymization irreversibility: computational, information-theoretic, and physics-guaranteed. We show these form a strict hierarchy (Section~\ref{sec:definitions}).

\item \textbf{Impossibility result.} We prove that no anonymization system whose randomness derives from a classical PRNG can achieve physics-guaranteed irreversibility (Theorem~\ref{thm:prng_impossible}).

\item \textbf{Construction.} We specify the \textsc{QRNG-OTP-Destroy} protocol and prove it achieves physics-guaranteed irreversibility under the Born rule assumption (Theorem~\ref{thm:qrng_secure}).

\item \textbf{Implementation.} We report on a production implementation with 10 progressive anonymization levels (L1 through L10), where L10 implements \textsc{QRNG-OTP-Destroy}. The system draws quantum entropy from IBM Quantum 156-qubit processors, Rigetti Ankaa, and qBraid, with automatic failover and provenance tracking.

\item \textbf{Regulatory analysis.} We argue that L10 output satisfies GDPR Recital~26's definition of anonymous information under a strictly stronger guarantee than any prior system, and that the quantum provenance log provides auditable evidence for data protection authorities.
\end{enumerate}

The remainder of this paper is organized as follows. Section~\ref{sec:background} reviews quantum measurement, classical PRNGs, and existing anonymization techniques. Section~\ref{sec:threat} defines the threat model. Section~\ref{sec:definitions} formalizes the three irreversibility tiers and proves the hierarchy. Section~\ref{sec:protocol} specifies the \textsc{QRNG-OTP-Destroy} protocol with its security proof. Section~\ref{sec:implementation} describes the implementation. Section~\ref{sec:comparison} presents a systematic comparison against existing tools. Section~\ref{sec:related} surveys related work. Section~\ref{sec:limitations} discusses limitations. Section~\ref{sec:conclusion} concludes.

%% ====================================================================
\section{Background}
\label{sec:background}
%% ====================================================================

\subsection{Quantum Measurement and the Born Rule}
\label{subsec:born}

A qubit is a two-level quantum system described by a state vector in a two-dimensional Hilbert space $\mathcal{H} = \mathbb{C}^2$. The computational basis states are $|0\rangle$ and $|1\rangle$. An arbitrary pure state is written
\begin{equation}
|\psi\rangle = \alpha|0\rangle + \beta|1\rangle, \quad |\alpha|^2 + |\beta|^2 = 1,
\label{eq:qubit}
\end{equation}
where $\alpha, \beta \in \mathbb{C}$ are probability amplitudes. When $|\psi\rangle$ is measured in the computational basis, the probability of each outcome is given by the Born rule:
\begin{equation}
P(0) = |\alpha|^2, \qquad P(1) = |\beta|^2.
\label{eq:born}
\end{equation}
For a balanced superposition ($\alpha = \beta = 1/\sqrt{2}$), each outcome occurs with probability exactly $1/2$. The outcome is not merely unpredictable in practice; it is fundamentally indeterminate prior to measurement. No hidden state predetermines the result.

For $N$ qubits each prepared in $|{+}\rangle = \frac{1}{\sqrt{2}}(|0\rangle + |1\rangle)$ and measured independently, the joint outcome is a uniformly random $N$-bit string. The min-entropy of this source is exactly $N$ bits:
\begin{equation}
H_\infty(X_1, X_2, \ldots, X_N) = N.
\label{eq:minentropy}
\end{equation}
No classical source achieves $H_\infty = N$ without a seed at least $N$ bits long, because a deterministic process cannot produce more entropy than its input contains.

\subsection{Bell's Theorem and Experimental Verification}
\label{subsec:bell}

The claim that quantum measurement outcomes are fundamentally non-deterministic requires ruling out the possibility that a hidden variable, not described by quantum mechanics, predetermines the result. Bell's theorem~\cite{bell1964epr} provides the test. Consider two entangled particles shared between distant parties Alice and Bob. Each party chooses a measurement setting ($a$ or $a'$ for Alice, $b$ or $b'$ for Bob) and records a binary outcome ($\pm 1$). Bell defined the quantity
\begin{equation}
S = E(a, b) - E(a, b') + E(a', b) + E(a', b')
\label{eq:chsh}
\end{equation}
where $E(a, b) = \langle A_a B_b \rangle$ is the correlation between Alice's and Bob's outcomes. Any local hidden variable (LHV) theory satisfies the CHSH inequality~$|S| \leq 2$. Quantum mechanics predicts a maximum violation of $|S| = 2\sqrt{2} \approx 2.828$.

Aspect, Dalibard, and Roger~\cite{aspect1982epr} provided the first experimental confirmation by measuring polarization correlations of entangled photon pairs and observing $|S| = 2.70 \pm 0.05$, exceeding the LHV bound by over five standard deviations. Subsequent experiments closed individual loopholes (detection efficiency, locality, freedom of choice). Hensen et al.~\cite{hensen2015loophole} closed all three loopholes simultaneously in a single experiment using entangled electron spins separated by 1.3~km, establishing that LHV models are ruled out at the $p < 3.9 \times 10^{-31}$ level. The 2022 Nobel Prize in Physics, awarded to Aspect, Clauser, and Zeilinger, recognized this line of work.

The consequence for our purposes is direct: random bits produced by measuring qubits in balanced superposition are not generated by any deterministic process. There is no seed. There is no internal state that, if captured, would allow reconstruction of the measurement outcomes. This is a physical fact, not a computational assumption.

\subsection{Limitations of Classical Pseudo-Random Number Generators}
\label{subsec:prng}

A CSPRNG such as ChaCha20 or AES-CTR-DRBG is a deterministic function $G\colon \mathcal{S} \to \{0,1\}^n$ that expands a short seed $s \in \mathcal{S}$ into a long output stream. The security property is computational indistinguishability: for any probabilistic polynomial-time distinguisher $\mathcal{D}$,
\begin{equation}
\bigl|\Pr[\mathcal{D}(G(s)) = 1] - \Pr[\mathcal{D}(U_n) = 1]\bigr| \leq \mathrm{negl}(\lambda)
\label{eq:csprng_security}
\end{equation}
where $U_n$ denotes the uniform distribution over $\{0,1\}^n$ and $\lambda$ is the security parameter. This guarantee is conditional on two assumptions: (1)~the seed remains secret, and (2)~no efficient algorithm breaks the underlying primitive (e.g., the AES block cipher or the ChaCha20 stream cipher).

The seed, however, is a physical object. It resides in RAM, in kernel entropy pools, in hardware state. Table~\ref{tab:prng_attacks} summarizes known attack vectors for seed extraction.

\begin{table}[t]
\caption{Known Attack Vectors for CSPRNG Seed Extraction}
\label{tab:prng_attacks}
\begin{center}
\begin{tabular}{@{}lp{4.8cm}@{}}
\toprule
\textbf{Attack Class} & \textbf{Mechanism} \\
\midrule
Memory forensics & Direct read of process memory or kernel entropy pool via \texttt{/proc/pid/mem} or memory dump \\
Cold boot attack & DRAM remanence after power loss; seed bits persist for seconds to minutes \\
Spectre/Meltdown & Speculative execution leaks PRNG state across privilege boundaries \\
Insider access & System administrator reads PRNG state from live process \\
Core dump capture & Seed included in crash dump written to disk \\
VM introspection & Hypervisor reads guest memory containing seed \\
\bottomrule
\end{tabular}
\end{center}
\end{table}

If the seed is captured through any of these vectors, the entire output stream is deterministically recoverable. For anonymization, this means that an adversary who obtains the PRNG state at the time of anonymization can reverse the transformation completely.

We emphasize the distinction between the CSPRNG's computational security guarantee and the physical security of the seed. The former is a mathematical property; the latter is an operational property that depends on the hardware, operating system, and deployment environment. An anonymization system whose irreversibility relies on CSPRNG security is only as strong as the weakest link in the seed's physical protection chain. No amount of algorithmic sophistication in the CSPRNG design can compensate for seed exposure.

\subsection{Classical Anonymization Techniques}
\label{subsec:classical_anon}

We review the principal anonymization techniques against which our system is compared.

\textbf{$K$-anonymity.} Sweeney~\cite{sweeney2002kanon} defined a dataset as $k$-anonymous if every record is indistinguishable from at least $k{-}1$ other records on quasi-identifier (QI) attributes. Formally, let $\mathrm{QI}(r)$ denote the quasi-identifier projection of record~$r$. A dataset $D$ satisfies $k$-anonymity if for every $r \in D$, $|\{r' \in D : \mathrm{QI}(r') = \mathrm{QI}(r)\}| \geq k$. The technique is vulnerable to homogeneity attacks (all records in an equivalence class share the same sensitive value) and background knowledge attacks.

\textbf{$\ell$-Diversity.} Machanavajjhala et al.~\cite{machanavajjhala2007ldiv} extended $k$-anonymity by requiring that each equivalence class contains at least $\ell$ ``well-represented'' values of each sensitive attribute, addressing the homogeneity vulnerability.

\textbf{$t$-Closeness.} Li et al.~\cite{li2007tcloseness} further strengthened the model by requiring that the distribution of sensitive attributes within each equivalence class is within distance $t$ (measured by Earth Mover's Distance) of the global distribution.

\textbf{Differential privacy.} Dwork et al.~\cite{dwork2006dp} introduced a fundamentally different approach. A randomized mechanism $\mathcal{M}$ satisfies $\epsilon$-differential privacy if for all datasets $D_1, D_2$ differing in one record and all measurable sets $S$:
\begin{equation}
\Pr[\mathcal{M}(D_1) \in S] \leq e^\epsilon \cdot \Pr[\mathcal{M}(D_2) \in S].
\label{eq:dp}
\end{equation}
The standard implementation adds Laplace noise with scale $\Delta f / \epsilon$, where $\Delta f$ is the sensitivity of the query function~$f$. Dwork and Roth~\cite{dwork2014algfound} provide the definitive treatment, covering composition theorems, the exponential mechanism, and connections to learning theory. Differential privacy is a property of the mechanism, not of the dataset: it bounds the information leakage per query, not the identifiability of the underlying records.

All of these techniques, when implemented in practice, draw their randomness from CSPRNGs. The $k$-anonymity family uses deterministic transformations (generalization, suppression) and does not require randomness at all, but the privacy guarantee is syntactic and has known composition vulnerabilities. Differential privacy requires high-quality randomness for noise generation, and its guarantee degrades if the noise can be reconstructed by an adversary who captures the CSPRNG seed. Both families therefore inherit the seed-capture vulnerability described above, either directly (DP) or through the weakness of deterministic structural guarantees ($k$-anonymity family).

\subsection{Regulatory Framework: GDPR and DORA}
\label{subsec:gdpr}

Two EU regulations define the legal context for our work.

\textbf{GDPR Recital~26.} The GDPR states that the principles of data protection ``should not apply to anonymous information, namely information which does not relate to an identified or identifiable natural person or to personal data rendered anonymous in such a manner that the data subject is not or no longer identifiable.'' The recital specifies that identifiability should be assessed considering ``all the means reasonably likely to be used'' for re-identification. The phrase ``reasonably likely'' has been interpreted by data protection authorities as requiring consideration of both current and foreseeable future capabilities. An anonymization method whose irreversibility depends on computational assumptions that may weaken over time faces an increasingly difficult argument under this standard.

The practical consequence is a two-tier classification with distinct regulatory burdens:

\begin{itemize}
\item \textbf{Anonymous data} (Recital~26): outside GDPR scope entirely. No consent requirements, no data subject access requests, no retention limits, no cross-border transfer restrictions.
\item \textbf{Pseudonymous data} (Article~4(5)): remains personal data. Full GDPR compliance required, including lawful basis for processing, data protection impact assessments, and right to erasure.
\end{itemize}

The economic difference between these two categories is substantial. An organization that can demonstrate true anonymization can freely share, archive, and process data. An organization limited to pseudonymization cannot.

\textbf{DORA Article~6.} The Digital Operational Resilience Act (DORA), in force for Norwegian financial institutions since July~2025, requires in Article~6.4 that ICT risk management frameworks include ``periodic cryptographic updates based on developments in cryptanalysis.'' This is the quantum-readiness clause: organizations must plan for the degradation of current cryptographic protections. Article~7 requires full cryptographic key lifecycle management. For anonymization systems that depend on CSPRNG security, DORA creates an ongoing obligation to re-assess whether the computational hardness assumption underlying the anonymization remains adequate. \textsc{QRNG-OTP-Destroy} eliminates this obligation: the physics-guaranteed irreversibility does not degrade with advances in cryptanalysis.

%% ====================================================================
\section{Threat Model}
\label{sec:threat}
%% ====================================================================

We consider four adversary classes, each strictly more powerful than the last. The goal of each adversary is to recover the original PII values from an anonymized dataset. Table~\ref{tab:adversaries} summarizes the four classes and their capabilities.

\begin{table*}[t]
\caption{Adversary Classes and Their Capabilities}
\label{tab:adversaries}
\begin{center}
\begin{tabular}{@{}llp{6cm}lp{4cm}@{}}
\toprule
\textbf{Class} & \textbf{Name} & \textbf{Capabilities} & \textbf{CSPRNG-Based} & \textbf{QRNG-OTP-Destroy} \\
\midrule
$\mathcal{A}_1$ & External, bounded compute & Anonymized dataset, knowledge of algorithm, current hardware & Secure & Secure \\
$\mathcal{A}_2$ & External, unbounded classical & Same as $\mathcal{A}_1$ plus unlimited classical compute. Brute-forces any seed space. Models P$=$NP scenario & \textbf{Broken} & Secure \\
$\mathcal{A}_3$ & External, quantum compute & Same as $\mathcal{A}_2$ plus fault-tolerant quantum computer. Runs Shor, Grover & \textbf{Degraded} & Secure \\
$\mathcal{A}_4$ & Insider, memory access & Memory snapshot of anonymization server at time of operation. Possesses PRNG state, intermediate buffers & \textbf{Broken} & Secure \\
\bottomrule
\end{tabular}
\end{center}
\end{table*}

\textbf{Adversary $\mathcal{A}_1$: External with bounded compute.} $\mathcal{A}_1$ possesses the anonymized dataset and knowledge of the anonymization algorithm but has no access to the anonymization server's memory or internal state. $\mathcal{A}_1$ has classical computational resources bounded by current hardware capabilities. This is the standard adversary model assumed by most anonymization systems. All techniques from Section~\ref{subsec:classical_anon} are designed to resist~$\mathcal{A}_1$.

\textbf{Adversary $\mathcal{A}_2$: External with unbounded classical compute.} $\mathcal{A}_2$ possesses the same inputs as $\mathcal{A}_1$ but has unlimited classical computational resources. $\mathcal{A}_2$ can brute-force any CSPRNG seed space, invert any hash function, and solve any problem in NP in polynomial time. This adversary models the long-term HNDL scenario. If P$=$NP, $\mathcal{A}_2$ is realistic. All CSPRNG-based anonymization methods are insecure against~$\mathcal{A}_2$.

\textbf{Adversary $\mathcal{A}_3$: External with quantum compute.} $\mathcal{A}_3$ has access to a universal fault-tolerant quantum computer in addition to unlimited classical compute. $\mathcal{A}_3$ can run Shor's algorithm, Grover's algorithm, and any quantum algorithm. Grover's algorithm quadratically speeds up brute-force search, reducing effective security by half in bit length.

\textbf{Adversary $\mathcal{A}_4$: Insider with memory access.} $\mathcal{A}_4$ is a system administrator or attacker who has obtained a memory snapshot of the anonymization server at the time of anonymization. $\mathcal{A}_4$ possesses the PRNG internal state, all intermediate buffers, and the OTP mapping (if it still exists in memory). Against $\mathcal{A}_4$, every CSPRNG-based anonymization method is completely broken: the PRNG state enables deterministic reconstruction of all ``random'' values. This is the adversary that motivates our work.

\textsc{QRNG-OTP-Destroy} resists all four adversary classes. Against $\mathcal{A}_4$ specifically, the defense rests on two properties: (1)~the quantum random bytes that generated the replacement tokens have no seed or state that $\mathcal{A}_4$ could capture, because the Born rule ensures the measurement outcomes are fundamentally non-deterministic; and (2)~the OTP mapping is destroyed via multi-pass overwrite before $\mathcal{A}_4$ can capture it. Even if $\mathcal{A}_4$ captures a snapshot after the mapping is destroyed, recovery requires determining which quantum measurement outcomes produced each token, which is information-theoretically impossible.

\textbf{Temporal window analysis.} The mapping exists in volatile memory only during the execution of the anonymization function (Step~2 through Step~4 of Algorithm~\ref{alg:qrng_otp}). For a dataset with 50,000 unique values, the mapping construction and substitution complete in under 1~second on commodity hardware (Section~\ref{subsec:perf}). An adversary $\mathcal{A}_4$ who captures a memory snapshot during this window obtains the mapping and can trivially invert the anonymization. This is true for any anonymization system, classical or quantum, that maintains an intermediate mapping. The QRNG contribution is that after destruction, no path to recovery exists. For classical systems, even after mapping destruction, the CSPRNG seed may persist elsewhere in memory, enabling full reconstruction. For \textsc{QRNG-OTP-Destroy}, no seed exists anywhere in the system at any point during execution.

The hardware enclave variant (Section~\ref{subsec:destruction}) further reduces the temporal window by isolating the mapping in enclave-protected memory, making it resistant to memory forensics even during execution.

%% ====================================================================
\section{Formal Definitions}
\label{sec:definitions}
%% ====================================================================

We formalize three tiers of anonymization irreversibility. Let $D$ be a dataset containing PII, let $A$ be an anonymization function, and let $D' = A(D)$ be the anonymized output. Let $\mathcal{A}$ denote an adversary attempting to recover $D$ from~$D'$.

\begin{definition}[Computational Irreversibility]
\label{def:comp}
An anonymization function $A$ is \emph{computationally irreversible} if, for every probabilistic polynomial-time adversary $\mathcal{A}$, the probability that $\mathcal{A}$ recovers any record of $D$ from $D'$ is negligible in the security parameter~$\lambda$:
\begin{equation}
\Pr[\mathcal{A}(D', 1^\lambda) \to d_i \in D] \leq \mathrm{negl}(\lambda).
\label{eq:comp_irrev}
\end{equation}
\end{definition}

This is the standard guarantee provided by CSPRNG-based anonymization. It holds under the assumption that the CSPRNG is secure, which requires the seed to remain secret and that no polynomial-time algorithm can distinguish the CSPRNG output from random.

\begin{definition}[Information-Theoretic Irreversibility]
\label{def:it}
An anonymization function $A$ is \emph{information-theoretically irreversible} if, for every adversary $\mathcal{A}$ with unbounded computational resources, the probability that $\mathcal{A}$ recovers any record of $D$ from $D'$ is bounded by:
\begin{equation}
\Pr[\mathcal{A}(D') \to d_i \in D] \leq 2^{-n}
\label{eq:it_irrev}
\end{equation}
where $n$ is the bit length of the replacement token.
\end{definition}

This guarantee is independent of computational assumptions. It requires that the replacement tokens are drawn from a distribution that is statistically independent of the original values and that no side information links the two.

\begin{definition}[Physics-Guaranteed Irreversibility]
\label{def:phys}
An anonymization function $A$ is \emph{physics-guaranteed irreversible} if its information-theoretic irreversibility (Definition~\ref{def:it}) holds not as a consequence of any mathematical assumption but as a consequence of a physical law. Specifically, the randomness source used by $A$ must satisfy the following: there exists no physical state, hidden variable, or prior information, in any physical theory consistent with observed experimental results, that determines the random values used in the anonymization transformation.
\end{definition}

Physics-guaranteed irreversibility is strictly stronger than information-theoretic irreversibility. The latter can, in principle, be achieved by a hypothetical perfect random number generator; the former requires that the randomness is certified by a physical theory with experimental support. The Born rule, validated to extraordinary precision by the loophole-free Bell test experiments~\cite{hensen2015loophole}, provides this certification.

\begin{lemma}[Strict Hierarchy]
\label{lem:hierarchy}
The three tiers form a strict hierarchy: physics-guaranteed irreversibility implies information-theoretic irreversibility implies computational irreversibility. The converses do not hold.
\end{lemma}

\begin{proof}
(i)~Physics-guaranteed $\Rightarrow$ information-theoretic: Definition~\ref{def:phys} requires that information-theoretic irreversibility (Definition~\ref{def:it}) holds as a consequence of physical law. This directly implies Definition~\ref{def:it}.

(ii)~Information-theoretic $\Rightarrow$ computational: If the bound in~\eqref{eq:it_irrev} holds for all adversaries (unbounded), it holds a fortiori for polynomial-time adversaries. For $n \geq \lambda$, the bound $2^{-n} \leq 2^{-\lambda}$ is negligible.

(iii)~Computational $\not\Rightarrow$ information-theoretic: A CSPRNG-based anonymization system is computationally irreversible (assuming the CSPRNG is secure) but not information-theoretically irreversible, because an unbounded adversary can enumerate the seed space $\mathcal{S}$ and reconstruct the mapping.

(iv)~Information-theoretic $\not\Rightarrow$ physics-guaranteed: Consider a hypothetical perfect TRNG that is information-theoretically random but whose physical mechanism is not certified by any tested physical theory (e.g., a thermal noise source whose randomness follows from assumptions about the noise model rather than from a fundamental law with loophole-free experimental verification). Such a source achieves Definition~\ref{def:it} but not Definition~\ref{def:phys}, because the randomness guarantee rests on a model assumption rather than an experimentally verified physical law.
\end{proof}

\begin{theorem}[Classical PRNG Impossibility]
\label{thm:prng_impossible}
No anonymization system whose randomness source is a classical PRNG (deterministic function of a finite seed) achieves physics-guaranteed irreversibility.
\end{theorem}

\begin{proof}
Let $A$ be an anonymization system that uses a PRNG $G\colon \mathcal{S} \to \{0,1\}^n$ with seed space~$\mathcal{S}$. The seed $s \in \mathcal{S}$ is a physical object stored in memory. An adversary who captures $s$ can compute $G(s)$, reconstruct every random value used in the anonymization, and invert $A$ to recover~$D$. The existence of $s$ as a physical state that determines the random values directly contradicts the condition in Definition~\ref{def:phys} that no physical state determines the random values. Therefore $A$ does not achieve physics-guaranteed irreversibility.
\end{proof}

\begin{theorem}[QRNG-OTP-Destroy Security]
\label{thm:qrng_secure}
The \textsc{QRNG-OTP-Destroy} protocol (Section~\ref{sec:protocol}) achieves physics-guaranteed irreversibility under the assumption that quantum mechanics correctly describes measurement outcomes (the Born rule assumption).
\end{theorem}

\begin{proof}
The protocol uses replacement tokens generated by measuring qubits in balanced superposition. By the Born rule, each measurement outcome is an independent uniformly random bit with no deterministic antecedent. By Bell's theorem and its loophole-free experimental verification~\cite{hensen2015loophole}, no local hidden variable determines the outcome. The OTP mapping between original values and replacement tokens is constructed in volatile memory and destroyed via multi-pass overwrite before the protocol returns.

After destruction, recovery of $D$ from $D'$ requires determining the quantum measurement outcomes that produced each replacement token. Since these outcomes are governed by the Born rule and have no deterministic antecedent, the probability of correctly guessing the replacement for any single value is $2^{-128}$ (for 16-byte tokens). This bound holds for any adversary, regardless of computational resources, because it follows from a physical law rather than a computational assumption.

The Born rule assumption is not a computational hardness assumption. It is a statement about the physical world, verified by experiment to extraordinary precision. Its failure would require a revision of quantum mechanics, contradicting over a century of experimental confirmation. Therefore the irreversibility of \textsc{QRNG-OTP-Destroy} is physics-guaranteed per Definition~\ref{def:phys}.
\end{proof}

\begin{corollary}[Independence from P vs.\ NP]
\label{cor:pvsnp}
The security of \textsc{QRNG-OTP-Destroy} is independent of the resolution of the $\mathrm{P}$ vs.\ $\mathrm{NP}$ problem.
\end{corollary}

\begin{proof}
Theorem~\ref{thm:qrng_secure} does not invoke any computational hardness assumption. If $\mathrm{P} = \mathrm{NP}$ were established, CSPRNG seed recovery would become polynomial-time, hash pre-image computation would become efficient, and every CSPRNG-based anonymization system would become reversible. The Born rule is a statement about measurement outcomes in quantum mechanics, not about the complexity of any computational problem. A world in which $\mathrm{P} = \mathrm{NP}$ but quantum mechanics remains valid is a world in which \textsc{QRNG-OTP-Destroy} remains secure and every CSPRNG-based anonymization system is broken.
\end{proof}

%% ====================================================================
\section{The QRNG-OTP-Destroy Protocol}
\label{sec:protocol}
%% ====================================================================

\subsection{Protocol Specification}
\label{subsec:spec}

The protocol takes as input a dataset $D$ (a table with $m$ columns and $n$ rows) and produces an anonymized dataset $D'$ of the same schema. It proceeds in four steps.

\begin{algorithm}[t]
\caption{QRNG-OTP-Destroy}
\label{alg:qrng_otp}
\begin{algorithmic}[1]
\REQUIRE Dataset $D$ with $m$ columns, $n$ rows; quantum entropy pool~$\mathcal{P}$
\ENSURE Anonymized dataset $D'$; all mappings destroyed
\FOR{each column $C_j$, $1 \leq j \leq m$}
    \STATE $U_j \leftarrow \text{unique}(C_j)$ \COMMENT{Set of unique values}
    \STATE $M_j \leftarrow \emptyset$ \COMMENT{Mapping in volatile memory}
    \FOR{each $v_k \in U_j$}
        \STATE Read 16 bytes $b_k$ from $\mathcal{P}$ \COMMENT{128 QRNG bits}
        \STATE $t_k \leftarrow \textsc{Base62}(b_k)$ \COMMENT{16-char token}
        \STATE $M_j[v_k] \leftarrow t_k$
    \ENDFOR
    \FOR{each row $i$, $1 \leq i \leq n$}
        \STATE $D'[i,j] \leftarrow M_j[D[i,j]]$
    \ENDFOR
\ENDFOR
\FOR{each mapping $M_j$, $1 \leq j \leq m$}
    \STATE \textsc{Overwrite}($M_j$, $\texttt{0x00}$) \COMMENT{Pass 1: zeros}
    \STATE \textsc{Overwrite}($M_j$, $\texttt{0xFF}$) \COMMENT{Pass 2: ones}
    \STATE \textsc{Overwrite}($M_j$, $\mathcal{P}$) \COMMENT{Pass 3: QRNG bytes}
    \STATE \textsc{Release}($M_j$)
\ENDFOR
\RETURN $D'$
\end{algorithmic}
\end{algorithm}

\textbf{Step~1: Entropy Acquisition.} For each column $C_j$ ($1 \leq j \leq m$), let $U_j = \{v_1, v_2, \ldots, v_{u_j}\}$ be the set of unique values in column~$C_j$. The protocol reads $16 \cdot \sum_{j=1}^{m} u_j$ bytes from a quantum entropy source. Each 16-byte block is produced by measuring 128 qubits, each prepared in the balanced superposition state $|{+}\rangle = \frac{1}{\sqrt{2}}(|0\rangle + |1\rangle)$, in the computational basis. The measurement outcomes are concatenated to form the 16-byte block.

\textbf{Step~2: Mapping Construction.} For each column $C_j$, the protocol constructs a mapping $M_j\colon U_j \to T$ where $T$ is the set of 16-character alphanumeric strings. Each unique value $v_k \in U_j$ is assigned a replacement token $t_k$ by converting the corresponding 16-byte QRNG block to a string using modular selection from the character set $\{$\texttt{a}--\texttt{z}, \texttt{A}--\texttt{Z}, \texttt{0}--\texttt{9}$\}$. The mapping $M_j$ is stored in volatile memory (a hash table in process address space).

\textbf{Step~3: Substitution.} Every cell $D[i,j]$ is replaced by $M_j(D[i,j])$, producing the anonymized dataset~$D'$. Identical values within a column map to the same token (preserving referential integrity within a single anonymization run), but different runs produce different tokens (non-reproducibility across runs).

\textbf{Step~4: Mapping Destruction.} All mappings $M_1, \ldots, M_m$ are destroyed. The destruction procedure overwrites the memory region occupied by each mapping with three passes: (1)~all zeros, (2)~all ones, (3)~random bytes from the quantum entropy source, following the DoD~5220.22-M standard. After overwrite, the memory is released. In an enhanced embodiment, the mapping is constructed within a hardware security enclave (Intel SGX or ARM TrustZone), and destruction is performed by enclave teardown.

\textbf{Entropy consumption.} Let $u = \sum_{j=1}^{m} |U_j|$ be the total number of unique values across all columns. The protocol consumes exactly $16u$ bytes for the OTP tokens plus $16u$ bytes for the third overwrite pass (QRNG bytes), for a total of $32u$ bytes per anonymization operation. For the destruction passes, the zero and one overwrites do not consume entropy. The total QRNG consumption is therefore:
\begin{equation}
E_{\mathrm{total}} = 32 \cdot \sum_{j=1}^{m} |U_j| \;\;\text{bytes}.
\label{eq:entropy_consumption}
\end{equation}

Table~\ref{tab:entropy_budget} shows entropy requirements for representative dataset sizes.

\begin{table}[t]
\caption{Entropy Consumption for Representative Datasets}
\label{tab:entropy_budget}
\begin{center}
\begin{tabular}{@{}rrrc@{}}
\toprule
\textbf{Rows} & \textbf{Columns} & \textbf{Unique Values} & \textbf{QRNG Bytes} \\
\midrule
1,000 & 5 & ${\sim}$3,000 & 96~KB \\
10,000 & 10 & ${\sim}$50,000 & 1.6~MB \\
100,000 & 20 & ${\sim}$500,000 & 16~MB \\
1,000,000 & 50 & ${\sim}$5,000,000 & 160~MB \\
\bottomrule
\end{tabular}
\end{center}
\end{table}

The entropy budget scales linearly with the number of unique values, not with the total number of cells. Datasets with high redundancy (many repeated values) consume far less entropy than datasets with mostly unique values. At the largest scale (1M rows, 50 columns), the 160~MB requirement is non-trivial but feasible with a dedicated QRNG appliance producing entropy at rates of 1~Gbit/s or higher (commercially available from ID~Quantique and Quantinuum).

\subsection{Security Analysis}
\label{subsec:security}

We state the per-value security bound and prove it directly.

\begin{proposition}[Per-value recovery bound]
\label{prop:pervalue}
After protocol execution, no adversary $\mathcal{A}$ with arbitrary computational resources (classical or quantum) can recover $D[i,j]$ from $D'[i,j]$ with probability exceeding $2^{-128}$.
\end{proposition}

\begin{proof}
We consider the adversary's information after the protocol completes.

(a) $\mathcal{A}$ possesses $D'$ and full knowledge of the protocol.

(b) $\mathcal{A}$ does not possess any mapping~$M_j$, because all mappings have been destroyed. Even an adversary with physical access to the server's memory after Step~4 finds only overwritten bytes.

(c) To recover $D[i,j]$ from $D'[i,j]$, $\mathcal{A}$ must determine which original value was mapped to the token $D'[i,j]$. This requires either (i)~inverting the mapping, which requires possessing the mapping, which has been destroyed; or (ii)~determining which QRNG output was assigned to each original value, which requires determining the quantum measurement outcomes that produced each 16-byte block.

(d) By the Born rule~\eqref{eq:born}, each measurement outcome is an independent fair coin flip with no deterministic antecedent. By Bell's theorem~\cite{bell1964epr}, no local hidden variable determines the outcome. The 128-bit QRNG block assigned to any particular original value is therefore uniformly distributed over $\{0,1\}^{128}$, independently of the original value. The probability that $\mathcal{A}$ correctly guesses the mapping for a single value is~$2^{-128}$.

(e) This bound is information-theoretic (it holds for unbounded adversaries) and physics-guaranteed (it follows from the Born rule rather than any computational assumption).
\end{proof}

\subsection{The Mapping Destruction Requirement}
\label{subsec:destruction}

The protocol's security critically depends on the mapping's destruction. If the mapping persists in any form, whether in memory, on disk, in a backup, in a log file, or in a core dump, an adversary who obtains it can trivially invert the anonymization. The QRNG entropy source provides unconditional randomness, but the mapping is classical information that must be treated as a one-time secret.

This constraint is no different in kind from the key management requirements of any encryption system. The distinction is that once the mapping is destroyed, the irreversibility guarantee is unconditional: it does not depend on the secrecy of any ongoing key or state.

For production deployments, we recommend:
\begin{enumerate}
\item Constructing the mapping in a hardware security enclave (Intel SGX, ARM TrustZone).
\item Disabling core dumps during the anonymization call.
\item Running the anonymization process in a memory-locked address space (\texttt{mlock}).
\item Verifying the overwrite by reading back the memory region after destruction.
\end{enumerate}

\subsection{Security Under P = NP}
\label{subsec:pnp}

Classical anonymization methods derive irreversibility from computational hardness. If P$=$NP were established, the implications would cascade through every CSPRNG: seed recovery would become polynomial-time, hash pre-image computation would become efficient, and every noise-based or token-based anonymization system that used a CSPRNG would become reversible.

\textsc{QRNG-OTP-Destroy} is immune to this scenario. The protocol's security does not invoke any computational hardness assumption at any point. The Born rule is a statement about measurement outcomes in quantum mechanics, not about the complexity of any computational problem. This independence from complexity-theoretic assumptions is formalized in Corollary~\ref{cor:pvsnp}.

%% ====================================================================
\section{Implementation}
\label{sec:implementation}
%% ====================================================================

\subsection{System Architecture}
\label{subsec:arch}

We have implemented \textsc{QRNG-OTP-Destroy} as Level~10 (L10) within Zipminator, a post-quantum cryptography platform with 10 progressive anonymization levels. Table~\ref{tab:levels} summarizes the complete hierarchy. Levels~1 through~9 implement classical techniques with increasing privacy strength. Level~10 implements the full \textsc{QRNG-OTP-Destroy} protocol.

\begin{table}[t]
\caption{Zipminator Anonymization Levels}
\label{tab:levels}
\begin{center}
\begin{tabular}{@{}clll@{}}
\toprule
\textbf{Level} & \textbf{Technique} & \textbf{Entropy} & \textbf{Irreversibility} \\
\midrule
L1 & Regex masking & None & Structural \\
L2 & SHA-3 hashing & None & Computational \\
L3 & Format-preserving tokenization & CSPRNG & Computational \\
L4 & $k$-anonymity ($k{=}5$) & None & Syntactic \\
L5 & $\ell$-diversity ($\ell{=}3$) & None & Syntactic \\
L6 & $t$-closeness ($t{=}0.2$) & None & Syntactic \\
L7 & Quantum noise jitter & QRNG & Computational \\
L8 & Differential privacy ($\epsilon$) & CSPRNG & Computational \\
L9 & $k$-anon $+$ DP combined & CSPRNG & Computational \\
L10 & \textsc{QRNG-OTP-Destroy} & QRNG & \textbf{Physics} \\
\bottomrule
\end{tabular}
\end{center}
\end{table}

The hierarchy is designed so that each level provides strictly stronger privacy than the one below it, at the cost of reduced analytical utility. Levels~1--3 are reversible by design (the mapping or hash can be stored for future lookups). Levels~4--6 apply syntactic models that discard information but are vulnerable to composition and background knowledge attacks. Levels~7--9 add noise from various entropy sources. Level~10 is the only level that provides information-theoretic, physics-guaranteed irreversibility.

The implementation is structured in three layers:

\textbf{Quantum entropy layer.} A background harvester service executes quantum circuits on IBM Quantum processors (156-qubit Fez and Marrakesh systems) via the qBraid gateway API. Each circuit prepares $N$ qubits in the state $|{+}\rangle^{\otimes N}$ and measures all qubits in the computational basis, yielding $N$ random bits per shot. The measurement outcomes are written to a binary entropy pool file (\texttt{quantum\_entropy\_pool.bin}). The harvester supports multiple providers with automatic failover, arranged in the priority chain shown in Table~\ref{tab:providers}.

\begin{table}[t]
\caption{Entropy Provider Priority Chain}
\label{tab:providers}
\begin{center}
\begin{tabular}{@{}clcc@{}}
\toprule
\textbf{Priority} & \textbf{Provider} & \textbf{Type} & \textbf{Certified} \\
\midrule
1 & PoolProvider (pre-harvested) & QRNG & Yes \\
2 & qBraid gateway & QRNG & Yes \\
3 & IBM Quantum direct & QRNG & Yes \\
4 & Rigetti Ankaa & QRNG & Yes \\
5 & API-based QRNG & QRNG & Yes \\
6 & OS entropy (\texttt{/dev/urandom}) & CSPRNG & No \\
\bottomrule
\end{tabular}
\end{center}
\end{table}

When the OS fallback (priority~6) is used, the system marks the output as ``classically anonymized'' rather than ``quantum-certified,'' preserving the integrity of the physics-guaranteed claim. Each harvesting cycle appends approximately \SI{50}{\kilo\byte} of quantum random bytes. The pool supports concurrent reads with per-consumer offset tracking; consumed bytes are never re-read.

\textbf{Anonymization engine.} The \texttt{LevelAnonymizer} class in Python exposes the method \texttt{apply(df, level=10)}, which accepts a Pandas DataFrame and an anonymization level. For L10, the engine iterates over each column, reads 16 bytes from the entropy pool for each unique value, constructs the OTP mapping in a Python dictionary (volatile memory), performs the substitution, and discards the mapping when the function returns.

\textbf{Cryptographic core.} The underlying post-quantum infrastructure is implemented in Rust (ML-KEM-768 per NIST FIPS~203) with Python bindings via PyO3. While L10 anonymization does not use lattice-based cryptography directly, the entropy pool infrastructure, key management, and secure communication channels that transport entropy from quantum hardware to the anonymization engine are protected by post-quantum key encapsulation.

\subsection{The LevelAnonymizer API}
\label{subsec:api}

The public API is minimal by design:

\begin{verbatim}
from zipminator.anonymizer import LevelAnonymizer

anonymizer = LevelAnonymizer(
    entropy_pool_path="/path/to/pool.bin"
)
anonymized_df = anonymizer.apply(
    original_df, level=10
)
\end{verbatim}

The \texttt{apply} method performs the four protocol steps (entropy acquisition, mapping construction, substitution, mapping destruction) within a single synchronous call. The mapping exists only for the duration of the call and is not accessible to the caller. The method returns a new DataFrame with all values replaced.

Consistency is preserved within a single call: identical values in a column produce identical tokens. Across calls, the same original value produces different tokens because each call draws fresh QRNG bytes. This non-reproducibility is a security feature: it ensures that cross-run correlation is impossible.

\subsection{Multi-Provider Entropy with Provenance}
\label{subsec:provenance}

The entropy pool supports provenance tracking. Each harvesting cycle records a structured log entry (JSONL format) containing: the quantum hardware provider, processor identifier and qubit count, circuit type, timestamp, number of bytes harvested, and the pool offset range to which the bytes were written. A representative log entry is:

\begin{verbatim}
{"timestamp": "2026-03-23T14:30:00Z",
 "provider": "ibm_quantum",
 "processor": "ibm_fez",
 "qubits": 156,
 "circuit_type": "random_measurement",
 "bytes_harvested": 4096,
 "pool_offset": [1048576, 1052672],
 "certification": "born_rule"}
\end{verbatim}

This provenance log serves as auditable evidence for data protection authorities: it documents that the entropy used for a particular anonymization operation originated from quantum hardware, not from a classical PRNG. The combination of provenance log and anonymization timestamp allows an auditor to verify the full chain: quantum circuit execution $\to$ entropy pool $\to$ anonymization operation $\to$ output dataset.

The provider factory implements health monitoring and automatic failover. If the primary quantum provider returns an error or exceeds a configurable latency threshold, the system attempts the next provider in the priority chain (Table~\ref{tab:providers}). The OS fallback is used only as a last resort and is distinguished in both the provenance log (with \texttt{certification: ``csprng\_fallback''}) and the output metadata.

\subsection{Test Results}
\label{subsec:tests}

The anonymization engine has been validated with 109 tests:

\begin{itemize}
\item 64 unit tests covering all 10 anonymization levels, including edge cases (empty DataFrames, single-row datasets, columns with all identical values, mixed data types).
\item 45 integration tests verifying end-to-end behavior from entropy pool reads through substitution and output validation.
\end{itemize}

Key properties verified by the test suite: (1)~L10 output differs from input for every cell; (2)~identical input values within a column produce identical tokens within a single run; (3)~different runs produce different tokens for the same input; (4)~the mapping dictionary is not accessible after \texttt{apply} returns; (5)~the system falls back to OS entropy and marks the output accordingly when no quantum provider is available.

\subsection{Performance}
\label{subsec:perf}

Anonymization throughput is dominated by entropy pool I/O rather than by the substitution logic. Table~\ref{tab:performance} summarizes the performance characteristics.

\begin{table}[t]
\caption{Anonymization Performance (10,000 rows, 10 columns, ${\sim}$50,000 unique values)}
\label{tab:performance}
\begin{center}
\begin{tabular}{@{}lcc@{}}
\toprule
\textbf{Metric} & \textbf{L8 (DP)} & \textbf{L10 (QRNG-OTP)} \\
\midrule
Entropy consumed & ${\sim}$40~KB & ${\sim}$800~KB \\
Wall-clock time & ${\sim}$1.5~s & ${\sim}$2.0~s \\
Dominant cost & Pandas apply & Pool I/O \\
Network latency & None & None (pre-harvested) \\
\bottomrule
\end{tabular}
\end{center}
\end{table}

For the benchmark dataset (10,000 rows, 10 columns, approximately 50,000 unique values), L10 anonymization consumes approximately \SI{800}{\kilo\byte} of quantum entropy (16 bytes per unique value) and completes in under 2 seconds on commodity hardware. The entropy pool is pre-populated by the background harvester; the anonymization call itself performs sequential reads from the pool file with no network latency.

%% ====================================================================
\section{Systematic Comparison}
\label{sec:comparison}
%% ====================================================================

We compare Zipminator L10 against the principal open-source anonymization tools available as of early 2026. Our analysis focuses on four dimensions: (1)~the technique employed, (2)~the entropy source underlying any randomized operations, (3)~the basis on which irreversibility is claimed, and (4)~the implications for regulatory compliance under GDPR Recital~26.

\begin{table*}[t]
\caption{Comparison of Anonymization Tools}
\label{tab:comparison}
\begin{center}
\begin{tabular}{@{}llllcp{3.5cm}@{}}
\toprule
\textbf{Tool} & \textbf{Technique} & \textbf{Entropy Source} & \textbf{Irreversibility} & \textbf{P=NP} & \textbf{GDPR Recital 26} \\
\midrule
\textbf{Zipminator L10} & QRNG-OTP + destroy & IBM Quantum (156Q) & Physics (Born rule) & Secure & Full: information-theoretic anonymization \\
ARX~\cite{prasser2014arx} & $k$-anon, $\ell$-div, $t$-close, DP & Java \texttt{SecureRandom} & Computational & Broken & Partial: parameter-dependent \\
sdcMicro~\cite{templ2017sdc} & $k$-anon, microaggregation, PRAM & R Mersenne Twister & Computational & Broken & Partial: linkable structure \\
Google DP~\cite{wilson2020dpsql} & Laplace/Gaussian mechanism & \texttt{/dev/urandom} & Computational & Broken & Partial: noise addition only \\
Apple Local DP~\cite{apple2017dp} & Randomized response, CMS & Device CSPRNG & Computational & Broken & Partial: tunable $\epsilon$ \\
OpenDP & Laplace, Gaussian, exponential & Platform CSPRNG & Computational & Broken & Partial: composition-dependent \\
Amnesia & $k$-anon, generalization & Java CSPRNG & Computational & Broken & Partial: syntactic model \\
MS Presidio & PII detection + masking & Platform CSPRNG & Computational & Broken & Partial: operator-dependent \\
\bottomrule
\end{tabular}
\end{center}
\end{table*}

\subsection{Entropy Sources}
\label{subsec:entropy_sources}

Every tool in Table~\ref{tab:comparison} other than Zipminator L10 relies on classical pseudo-random number generators. ARX and Amnesia, being Java applications, use \texttt{java.security.SecureRandom}, which delegates to the platform CSPRNG. sdcMicro defaults to R's Mersenne Twister, a fast but non-cryptographic PRNG; users can substitute a CSPRNG, but the default configuration does not enforce this. Google's DP library and OpenDP use the operating system's CSPRNG for noise generation. Apple's local DP implementation uses the device's hardware-backed CSPRNG (Secure Enclave on iOS/macOS). Microsoft Presidio uses platform randomness for masking operations.

All of these generators are deterministic: given the internal state, every output is reproducible. The security guarantee is that recovering the state from the output is computationally infeasible. This is a reasonable assumption under current hardware, but it is an assumption, not a physical law.

\subsection{Irreversibility Analysis}
\label{subsec:irrev_analysis}

The tools fall into three categories based on their irreversibility claims.

\textbf{Syntactic methods} (ARX, sdcMicro, Amnesia) apply transformations such as generalization, suppression, and microaggregation. These are structurally irreversible in the sense that information is discarded during generalization (e.g., replacing an exact age with an age range). The degree of protection depends on parameter choices ($k$, $\ell$, $t$) and is vulnerable to composition attacks (multiple releases of the same dataset can be cross-referenced), homogeneity attacks (all records in an equivalence class share the same sensitive value), and background knowledge attacks (external information narrows the candidate set)~\cite{machanavajjhala2007ldiv, li2007tcloseness}. No syntactic method makes a claim about the entropy source because the transformation itself is deterministic.

\textbf{Noise-addition methods} (Google DP, Apple Local DP, OpenDP) add calibrated random noise to query outputs or individual records. The privacy guarantee is differential privacy with a specified~$\epsilon$. The noise is generated from a CSPRNG. If the CSPRNG state is compromised through any of the attack vectors in Table~\ref{tab:prng_attacks}, the noise values can be reconstructed and subtracted, reversing the anonymization entirely. The guarantee is therefore computational: it holds as long as the adversary cannot recover the CSPRNG state. Under the HNDL threat model, an adversary who captures the CSPRNG state today can reverse the DP guarantee at any future time.

\textbf{Zipminator L10} replaces every value with a quantum-random string and destroys the mapping. The irreversibility rests on two independent foundations: (1)~the replacement values are generated from quantum measurements with no deterministic seed, and (2)~the mapping between original and replacement values is destroyed via multi-pass overwrite. Reversing L10 requires either predicting quantum measurement outcomes (impossible by the Born rule) or recovering a destroyed mapping (impossible given secure erasure). This is the only tool in the comparison that achieves information-theoretic irreversibility.

Table~\ref{tab:irrev_summary} summarizes the irreversibility guarantees across the three categories.

\begin{table}[t]
\caption{Summary of Irreversibility Guarantees}
\label{tab:irrev_summary}
\begin{center}
\begin{tabular}{@{}lcccc@{}}
\toprule
\textbf{Category} & \textbf{$\mathcal{A}_1$} & \textbf{$\mathcal{A}_2$} & \textbf{$\mathcal{A}_3$} & \textbf{$\mathcal{A}_4$} \\
\midrule
Syntactic ($k$-anon) & Partial & Partial & Partial & N/A \\
Noise-addition (DP) & Secure & Broken & Degraded & Broken \\
QRNG-OTP-Destroy & Secure & Secure & Secure & Secure \\
\bottomrule
\end{tabular}
\end{center}
\end{table}

\subsection{Why This Has Not Been Done Before}
\label{subsec:gap}

The conceptual gap is disciplinary, not technological. QRNG hardware has been commercially available since 2001 (ID~Quantique) and cloud-accessible since at least 2019 (IBM Quantum Experience). The anonymization community treated randomness as a solved problem: CSPRNGs are ``good enough'' for noise generation, and research focus has been on tightening $\epsilon$ bounds, improving utility, and defending against composition attacks. The QRNG community focused on cryptographic key generation, not on anonymization.

The result is that no existing tool connects quantum entropy to the irreversibility argument for data anonymization. The closest work, discussed in Section~\ref{sec:related}, is the \emph{Nature Reviews Physics} perspective by Amer et al.~\cite{amer2025certified}, which identifies differential privacy as a promising application of certified randomness but does not discuss anonymization, mapping destruction, or the information-theoretic irreversibility argument.

\subsection{QRNG Retrofitting Feasibility}
\label{subsec:retrofit}

Could existing tools swap their CSPRNG for a QRNG? In principle, yes: replacing the entropy source is an engineering task. This alone would not achieve information-theoretic irreversibility. The tools would also need to (1)~implement a one-time pad mapping scheme rather than noise addition, (2)~implement secure mapping destruction, and (3)~restructure their anonymization pipeline to treat the mapping as a volatile, single-use artifact. These are architectural changes, not parameter swaps. The combination of QRNG sourcing, OTP mapping, and mapping destruction is the novel contribution.

We note that even a QRNG-enhanced differential privacy mechanism (replacing CSPRNG noise with QRNG noise) would not achieve physics-guaranteed irreversibility in our sense. The noise values, once generated, exist as classical data in memory. If the noise values are captured before they are discarded, the adversary can subtract them. The contribution of QRNG to DP is that the noise values cannot be \emph{predicted} from any seed, but they can still be \emph{observed} if the adversary has memory access. \textsc{QRNG-OTP-Destroy} addresses this through the mapping destruction step: the classical artifact (the mapping) is eliminated, and the quantum artifact (the measurement outcomes) never existed as a recoverable state.

\subsection{Regulatory Implications}
\label{subsec:regulatory}

The comparison in Table~\ref{tab:comparison} has direct regulatory implications under both GDPR and DORA.

Under GDPR Recital~26, an anonymization method must render data subjects ``not or no longer identifiable'' considering ``all the means reasonably likely to be used.'' For CSPRNG-based methods, the argument that re-identification is not ``reasonably likely'' depends on the continued secrecy of the CSPRNG seed and the continued hardness of seed recovery. As quantum computing advances and memory forensics tools become more capable, the ``reasonably likely'' threshold shifts. Data anonymized today with a CSPRNG-based method may not satisfy Recital~26 in a decade if the computational assumptions weaken.

For \textsc{QRNG-OTP-Destroy}, the Recital~26 argument does not depend on any computational assumption. Re-identification requires predicting quantum measurement outcomes, which is physically impossible. The ``reasonably likely'' test is satisfied permanently: no foreseeable advance in technology changes the physics of quantum measurement.

Under DORA Article~6.4, financial institutions must perform ``periodic cryptographic updates based on developments in cryptanalysis.'' For CSPRNG-based anonymization, this creates an ongoing compliance obligation: the institution must regularly re-assess whether the anonymization remains adequate in light of computational advances. For \textsc{QRNG-OTP-Destroy}, no such re-assessment is needed, because the guarantee is not derived from cryptanalysis-vulnerable assumptions. This reduces the operational compliance burden.

%% ====================================================================
\section{Related Work}
\label{sec:related}
%% ====================================================================

\subsection{Differential Privacy}
\label{subsec:dp}

Dwork et al.~\cite{dwork2006dp} introduced differential privacy as a formal framework for privacy-preserving data analysis, proving that calibrating Laplace noise to the sensitivity of a query function bounds the influence of any single individual's data on the output. The key insight is that privacy is a property of the \emph{mechanism}, not of the dataset: the same mechanism provides the same $\epsilon$ guarantee regardless of the underlying data distribution. Dwork and Roth~\cite{dwork2014algfound} provide the definitive treatment, covering composition theorems (privacy loss accumulates across multiple queries), the exponential mechanism (for discrete outputs), and connections to learning theory.

Wilson et al.~\cite{wilson2020dpsql} developed differentially private SQL with bounded user contribution, demonstrating practical deployment at Google scale. Apple~\cite{apple2017dp} implemented local differential privacy for telemetry collection, where noise is added on the user's device before data leaves the device.

These results establish the gold standard for statistical privacy guarantees. Two limitations are relevant to our work. First, differential privacy protects the output of computations over data; the original data persists in the data holder's custody. Second, all practical DP implementations use classical PRNGs to generate noise, inheriting the seed-recovery vulnerability discussed in Section~\ref{subsec:prng}. An adversary who captures the PRNG state can subtract the noise and recover exact query results.

\subsection{$k$-Anonymity and Extensions}
\label{subsec:kanon}

Sweeney~\cite{sweeney2002kanon} proposed $k$-anonymity as a syntactic privacy model requiring that each record be indistinguishable from at least $k{-}1$ others on quasi-identifier attributes. The model has known weaknesses. Machanavajjhala et al.~\cite{machanavajjhala2007ldiv} demonstrated vulnerability to homogeneity attacks (all records in an equivalence class share the same sensitive value, making the sensitive attribute trivially inferable) and background knowledge attacks (external information narrows the candidate set below~$k$). They proposed $\ell$-diversity, requiring at least $\ell$ distinct sensitive values per equivalence class. Li et al.~\cite{li2007tcloseness} introduced $t$-closeness, requiring the distribution of sensitive attributes within equivalence classes to be within Earth Mover's Distance $t$ of the global distribution, addressing the ``skewness attack'' where a non-uniform distribution within an equivalence class leaks information even under $\ell$-diversity.

These models are deterministic: they apply generalization and suppression without randomness. Their guarantees are syntactic (structural properties of the output table) rather than probabilistic. They do not claim information-theoretic guarantees and are known to be vulnerable to composition when multiple releases of the same dataset are available.

\subsection{QRNG in Cryptography}
\label{subsec:qrng_crypto}

Quantum random number generators have been deployed commercially for cryptographic key generation since 2004, when ID~Quantique (Geneva) brought the first commercial QKD system to market. ID~Quantique's Quantis product line (USB, PCIe, chip, and appliance form factors) has received NIST SP~800-90B Entropy Source Validation on the IID track, making it the first QRNG to achieve this certification. Quantinuum's Quantum Origin platform generates cryptographic keys from verified quantum randomness produced on trapped-ion processors; in 2025, Quantum Origin became the first software QRNG to achieve NIST SP~800-90B validation.

Both platforms focus exclusively on key generation and key management. The conceptual step from ``QRNG makes better keys'' to ``QRNG makes irreversible anonymization'' has not been taken in the commercial QRNG literature. This is the disciplinary gap our work fills.

\subsection{Certified Randomness}
\label{subsec:certified}

Amer et al.~\cite{amer2025certified}, published in \emph{Nature Reviews Physics}, present a comprehensive survey of applications of certified randomness generated by quantum computers. The authors identify four application domains: cryptography, computational advantage demonstrations, randomness expansion, and statistical privacy. For the privacy domain, they argue that certified randomness can provide ``everlasting privacy'' against unbounded adversaries when used in place of pseudorandom noise in differential privacy mechanisms.

This is the closest existing work to our contribution. Three distinctions are critical:
\begin{enumerate}
\item Amer et al.\ discuss DP noise generation (replacing CSPRNG noise with QRNG noise), not data anonymization (replacing data values with QRNG-derived tokens).
\item They do not propose mapping destruction. Without mapping destruction, QRNG-enhanced DP still requires protecting intermediate state.
\item They do not make the information-theoretic irreversibility argument for rendered-anonymous data under GDPR Recital~26.
\end{enumerate}

Their work confirms the novelty of our approach: the connection between certified quantum randomness and data anonymization is recognized as a research frontier, but the specific construction we present (QRNG OTP with mapping destruction yielding physics-guaranteed irreversibility) has not appeared in the literature.

\subsection{Information-Theoretic Security: Shannon's OTP}
\label{subsec:it_security}

Shannon~\cite{shannon1949secrecy} proved that the one-time pad achieves perfect secrecy: for a message $m$ encrypted with a key $k$ drawn uniformly at random from the same-length key space, the ciphertext $c = m \oplus k$ reveals no information about~$m$. Formally, $H(M | C) = H(M)$: the conditional entropy of the message given the ciphertext equals the a priori entropy of the message. This is information-theoretic security; it holds against computationally unbounded adversaries.

Our contribution applies the same class of guarantee to a different problem. Rather than securing a communication channel, we secure a data transformation. The OTP mapping in L10 functions as a one-time pad applied to data values rather than to ciphertext. The critical distinction from classical OTP encryption is twofold: (1)~the ``key'' (mapping) is derived from QRNG rather than a shared secret, ensuring no seed can be captured; and (2)~the key is destroyed rather than shared with a recipient, ensuring no copy persists after the anonymization completes. The combination of quantum randomness and key destruction yields a guarantee that is strictly stronger than Shannon's OTP: not only is the ``ciphertext'' (anonymized data) information-theoretically independent of the ``plaintext'' (original data), but the ``key'' (mapping) is physically irrecoverable.

\subsection{Quantum Differential Privacy}
\label{subsec:qdp}

A distinct line of research, initiated by Hirche, Rouz\'{e}, and Stilck Fran\c{c}a (arXiv:2202.10717), develops an information-theoretic framework for quantum differential privacy. Their work recasts DP as a quantum divergence and shows that the inherent noise of near-term quantum computers provides natural differential privacy for quantum computations. This addresses a fundamentally different problem: protecting quantum states during quantum computation. Our work protects classical data using quantum randomness as an entropy source for classical anonymization. The two lines of research are complementary but non-overlapping.

%% ====================================================================
\section{Discussion}
\label{sec:discussion}
%% ====================================================================

\subsection{The Privacy-Utility Spectrum}
\label{subsec:spectrum}

Anonymization systems operate on a spectrum from full utility (no privacy) to full privacy (no utility). At one extreme, releasing the raw dataset preserves all analytical value but provides no privacy protection. At the other extreme, L10 replaces every value with quantum-random tokens, providing physics-guaranteed privacy but destroying all analytical utility.

The 10-level hierarchy in Table~\ref{tab:levels} is designed to let practitioners choose their position on this spectrum. For a dataset intended for statistical analysis, L8 (differential privacy with tunable $\epsilon$) may be appropriate: the noise preserves aggregate statistics while bounding individual information leakage. For a dataset intended to be provably destroyed as data while retaining the table structure (e.g., to satisfy a GDPR Article~17 deletion request without disrupting downstream systems that reference row IDs), L10 is the appropriate choice.

The key insight is that L10 is not a replacement for differential privacy in analytical workflows. It is a new category of technique for a use case that DP was not designed to address: rendering data permanently, provably anonymous under the strongest possible guarantee.

\subsection{Assumptions and Their Scope}
\label{subsec:assumptions}

The security of \textsc{QRNG-OTP-Destroy} rests on two assumptions:

\begin{enumerate}
\item \textbf{The Born rule.} Quantum measurement outcomes for qubits prepared in $|{+}\rangle$ are uniformly random with min-entropy $H_\infty = 1$ bit per qubit. This has been experimentally verified via loophole-free Bell tests at significance levels exceeding $10^{-30}$~\cite{hensen2015loophole}. A violation would constitute a fundamental revision of quantum mechanics.

\item \textbf{Secure mapping destruction.} The multi-pass overwrite (or enclave teardown) successfully eliminates all copies of the mapping. This is an operational assumption about the hardware and software stack, not a mathematical assumption. It can be verified through formal methods (memory safety proofs), hardware attestation (enclave integrity verification), and physical testing (cold boot attack resistance measurement).
\end{enumerate}

The first assumption is, in our assessment, the strongest assumption available in science: it is supported by over a century of experimental confirmation and has never been contradicted. The second assumption is an engineering requirement, comparable in kind to the key management requirements of any cryptographic system. The distinction from CSPRNG-based systems is that \textsc{QRNG-OTP-Destroy} has only one operational assumption (mapping destruction), whereas CSPRNG-based systems have two (seed secrecy and CSPRNG security), and the failure of either is catastrophic.

\subsection{Implications for the HNDL Threat}
\label{subsec:hndl}

The harvest-now, decrypt-later threat is well understood for encrypted data: an adversary captures ciphertext today and decrypts it when quantum computers become powerful enough to break RSA or ECDH. The migration path is post-quantum cryptography (FIPS~203, 204, 205).

For anonymized data, the HNDL threat has received less attention. An adversary who captures both the anonymized dataset and the CSPRNG state (or who captures the CSPRNG state and waits for computational advances that make seed recovery efficient) can reverse the anonymization at a later time. The migration path for anonymized data has, until now, been undefined.

\textsc{QRNG-OTP-Destroy} provides this migration path. Data anonymized with L10 is immune to the HNDL threat because: (1)~there is no seed to capture; (2)~there is no mapping to store; and (3)~the guarantee does not degrade with computational advances. For organizations that have already anonymized data using CSPRNG-based methods and are concerned about HNDL, re-anonymizing with L10 provides a permanent upgrade.

\subsection{Comparison with Quantum Key Distribution}
\label{subsec:qkd_comparison}

Quantum key distribution (QKD) is the best-known application of quantum mechanics to information security. QKD allows two parties to establish a shared secret key whose secrecy is guaranteed by quantum mechanics (specifically, the no-cloning theorem and the disturbance caused by eavesdropping on quantum channels). Our work shares the foundational principle, that quantum mechanics provides guarantees beyond computational hardness, but applies it to a different problem.

Table~\ref{tab:qkd_vs_l10} summarizes the comparison. The key difference is that QKD produces a shared secret (the key exists on both sides and must be protected), while L10 produces an anonymous dataset with no persisting secret (the mapping is destroyed). QKD requires a quantum channel between communicating parties; L10 requires only a quantum random number source, which can be pre-populated offline.

\begin{table}[t]
\caption{Comparison: QKD vs.\ QRNG-OTP-Destroy}
\label{tab:qkd_vs_l10}
\begin{center}
\begin{tabular}{@{}lll@{}}
\toprule
\textbf{Property} & \textbf{QKD} & \textbf{L10} \\
\midrule
Goal & Secure key exchange & Irreversible anonymization \\
Quantum resource & Entangled photons & Superposition measurements \\
Output & Shared secret key & Anonymous dataset \\
Persisting secret & Yes (the key) & No (mapping destroyed) \\
Quantum channel & Required (real-time) & Not required (offline) \\
Guarantee basis & No-cloning theorem & Born rule \\
Computational assumption & None & None \\
\bottomrule
\end{tabular}
\end{center}
\end{table}

This comparison illustrates an underexplored design space: applying quantum-mechanical guarantees to data privacy problems beyond key exchange. QKD secures communication; \textsc{QRNG-OTP-Destroy} secures data transformation. Other applications in this space, such as quantum-certified random beacons for auditable lotteries, quantum-enhanced secret sharing, and quantum-sourced noise for differential privacy, represent fertile ground for future work.

%% ====================================================================
\section{Limitations}
\label{sec:limitations}
%% ====================================================================

We identify five limitations of the L10 quantum anonymization approach.

\textbf{QRNG availability and cost.} L10 requires quantum random bytes produced by measuring qubits in superposition. As of 2026, this requires access to cloud-based quantum computing hardware (IBM Quantum, Rigetti, or equivalent) or dedicated QRNG appliances (ID~Quantique, Quantinuum). Cloud quantum access carries per-job costs and latency. QRNG appliances cost thousands to tens of thousands of dollars. Our entropy pool architecture mitigates latency by pre-harvesting quantum random bytes in a background daemon, but the pool is finite and must be replenished. Organizations without quantum hardware access fall back to OS-level entropy (\texttt{/dev/urandom}), which provides computational security only. The system marks such operations as ``classically anonymized'' rather than ``quantum-certified.''

\textbf{Statistical inference attacks.} QRNG protects against seed-recovery attacks: no adversary can reconstruct the replacement values by recovering a seed that does not exist. QRNG does not protect against statistical inference on the anonymized dataset itself. If the anonymized data preserves structural properties of the original (e.g., frequency distributions, correlations between columns), an adversary with background knowledge may infer sensitive attributes without reversing the OTP. L10's full-replacement strategy (every value is replaced) eliminates direct linkage, but the schema and row structure are preserved. For datasets where structural properties are sensitive, additional measures (row shuffling, synthetic data generation) may be needed.

\textbf{Utility destruction.} L10 is a maximum-privacy, zero-utility transformation. Every original value is replaced with a quantum-random string; the anonymized dataset cannot be used for statistical analysis, machine learning, or any computation that depends on the semantic content of the data. This is by design: L10 targets the use case where the goal is to render data provably anonymous for regulatory purposes (e.g., satisfying a data deletion request under GDPR Article~17 without destroying the dataset structure). For use cases requiring analytical utility, Zipminator offers lower anonymization levels (L1 through L9) with configurable privacy-utility tradeoffs.

\textbf{Pool exhaustion.} A single anonymization operation on a moderately sized dataset (10,000 rows, 10 columns, ${\sim}$100,000 unique values) consumes approximately \SI{1.6}{\mega\byte} of quantum entropy. The entropy pool must be replenished by executing quantum circuits on hardware, which is rate-limited by processor availability and queue times. Under heavy load, pool exhaustion is possible. Our implementation monitors pool health, triggers refill when remaining bytes fall below a configurable threshold, and refuses to perform quantum-certified anonymization when the pool is insufficient.

\textbf{Trust in quantum hardware vendors.} The Born rule guarantee assumes that the quantum hardware faithfully prepares qubits in superposition and that measurement outcomes are not manipulated. In practice, this means trusting that IBM, Rigetti, or the QRNG appliance vendor has not introduced a backdoor that makes measurement outcomes deterministic (e.g., by secretly reading a classical seed rather than performing genuine quantum measurement).

This trust assumption can be partially mitigated through three mechanisms:

\begin{enumerate}
\item \textbf{Multi-provider aggregation.} Our implementation draws entropy from multiple independent quantum hardware providers (Table~\ref{tab:providers}). Compromising the randomness would require collusion among all providers.

\item \textbf{Statistical testing.} The harvested entropy undergoes standard randomness tests (NIST SP~800-22 test suite) before being written to the pool. Biased or patterned outputs are rejected.

\item \textbf{Certified randomness protocols.} Device-independent randomness certification, as surveyed by Amer et al.~\cite{amer2025certified}, uses Bell inequality violations between two separated devices to verify that outputs are genuinely quantum without trusting the device implementation. Integrating such protocols into the entropy harvester is future work.
\end{enumerate}

\textbf{Formal verification gap.} Our mapping destruction implementation uses Python's \texttt{ctypes.memset} to overwrite memory. Python's garbage collector and memory allocator may retain copies of objects in internal pools or free lists. While our multi-pass overwrite targets the primary memory region, we cannot guarantee with mathematical certainty that no copy of any mapping entry persists in interpreter-managed memory. A formally verified implementation in Rust or C, compiled with memory-safe semantics and executed within a hardware enclave, would close this gap. This is planned for a future release.

\textbf{Column structure preservation.} L10 replaces all values but preserves the dataset's schema (column names, column count, row count). An adversary who knows the schema and has access to auxiliary information about the original dataset's structure (e.g., the number of unique values per column) may be able to infer column semantics. For example, a column with exactly 2 unique tokens likely corresponds to a binary attribute. This structural leakage is inherent to any anonymization method that preserves schema. For datasets where schema information is sensitive, column names should be anonymized separately, and dummy columns or rows may be added to obscure the structure.

%% ====================================================================
\section{Conclusion}
\label{sec:conclusion}
%% ====================================================================

We have presented the first anonymization system where irreversibility is guaranteed by quantum mechanics rather than computational hardness assumptions. The construction is direct: replace every data value with a quantum-random string generated from Born-rule measurements, then destroy the mapping. The security argument is equally direct: reversing the transformation requires either predicting quantum measurement outcomes (impossible by the Born rule and Bell's theorem~\cite{bell1964epr, hensen2015loophole}) or recovering a mapping that has been overwritten (impossible given secure erasure). No seed exists. No replay is possible. The guarantee holds against adversaries with unbounded computational resources, including universal quantum computers, and is independent of the resolution of the P vs.\ NP problem.

This construction fills a gap that has persisted despite the simultaneous availability of QRNG hardware and mature anonymization techniques. The gap was disciplinary, not technological: the QRNG community focused on key generation while the anonymization community treated classical randomness as sufficient. The closest prior work~\cite{amer2025certified} identifies certified randomness as promising for differential privacy but does not address anonymization, mapping destruction, or the GDPR Recital~26 compliance argument.

Data rendered anonymous by L10 satisfies GDPR Recital~26's definition of anonymous information, where the data subject is ``not or no longer identifiable.'' Unlike pseudonymized data, which remains subject to GDPR, L10 output falls outside the scope of data protection regulation entirely. For organizations subject to DORA Article~6, which requires cryptographic updates in response to advances in cryptanalysis, L10 provides a guarantee that does not degrade with advances in computing.

A provisional patent application has been filed (Norwegian Industrial Property Office, March 2026) covering the method and system for irreversible data anonymization using QRNG with physics-guaranteed non-reversibility. The implementation is open-source (Apache-2.0) and available as a Python SDK, REST API, and CLI tool. The anonymization engine, entropy pool, and multi-provider QRNG infrastructure are tested with 429 Python tests and 441 Rust tests.

Future work proceeds along five axes:

\begin{enumerate}
\item \textbf{Formal verification of mapping destruction.} We plan to formally verify, using tools such as Frama-C or CBMC, that the multi-pass overwrite procedure leaves no residual mapping data in any addressable memory location.

\item \textbf{Hardware security integration.} Production deployment should place the mapping inside a hardware security enclave (Intel SGX, ARM TrustZone, or AWS Nitro Enclaves) to eliminate the temporal window during which the mapping exists in general-purpose memory.

\item \textbf{Device-independent randomness certification.} Our current implementation trusts the quantum hardware vendor to produce genuine quantum randomness. Integrating device-independent certification protocols, which verify randomness quality using Bell inequality violations between two separated devices, would reduce this trust requirement.

\item \textbf{Hybrid schemes with analytical utility.} L10 is zero-utility by design. A natural extension is a hybrid scheme that applies QRNG-OTP to direct identifiers and sensitive attributes while preserving the statistical structure of quasi-identifiers through $k$-anonymity or DP, yielding a dataset that is partially analyzable and partially physics-guaranteed anonymous.

\item \textbf{Standardization engagement.} We intend to propose quantum-certified anonymization as a recognized privacy-enhancing technology to NIST, ENISA, and the Article~29 Working Party (now EDPB), with the goal of establishing QRNG-OTP-Destroy as a reference technique for Recital~26 compliance.
\end{enumerate}

%% ====================================================================
\begin{thebibliography}{14}

\bibitem{bell1964epr}
J.~S. Bell, ``On the Einstein Podolsky Rosen paradox,'' \emph{Physics Physique Fizika}, vol.~1, no.~3, pp.~195--200, 1964.

\bibitem{aspect1982epr}
A.~Aspect, J.~Dalibard, and G.~Roger, ``Experimental realization of Einstein-Podolsky-Rosen-Bohm Gedankenexperiment: A new violation of Bell's inequalities,'' \emph{Physical Review Letters}, vol.~49, no.~2, pp.~91--94, 1982.

\bibitem{hensen2015loophole}
B.~Hensen, H.~Bernien, A.~E. Dr\'{e}au, \emph{et al.}, ``Loophole-free Bell inequality violation using electron spins separated by 1.3 kilometres,'' \emph{Nature}, vol.~526, pp.~682--686, 2015.

\bibitem{shannon1949secrecy}
C.~E. Shannon, ``Communication theory of secrecy systems,'' \emph{Bell System Technical Journal}, vol.~28, no.~4, pp.~656--715, 1949.

\bibitem{sweeney2002kanon}
L.~Sweeney, ``$k$-Anonymity: A model for protecting privacy,'' \emph{International Journal of Uncertainty, Fuzziness and Knowledge-Based Systems}, vol.~10, no.~5, pp.~557--570, 2002.

\bibitem{dwork2006dp}
C.~Dwork, F.~McSherry, K.~Nissim, and A.~Smith, ``Calibrating noise to sensitivity in private data analysis,'' in \emph{Proc.\ Theory of Cryptography Conf.\ (TCC)}, ser.\ Lecture Notes in Computer Science, vol.~3876.\hskip 1em plus 0.5em minus 0.4em Springer, 2006, pp.~265--284.

\bibitem{machanavajjhala2007ldiv}
A.~Machanavajjhala, D.~Kifer, J.~Gehrke, and M.~Venkitasubramaniam, ``$\ell$-Diversity: Privacy beyond $k$-anonymity,'' \emph{ACM Transactions on Knowledge Discovery from Data}, vol.~1, no.~1, Article~3, 2007.

\bibitem{li2007tcloseness}
N.~Li, T.~Li, and S.~Venkatasubramanian, ``$t$-Closeness: Privacy beyond $k$-anonymity and $\ell$-diversity,'' in \emph{Proc.\ IEEE 23rd Int.\ Conf.\ Data Engineering (ICDE)}, 2007.

\bibitem{dwork2014algfound}
C.~Dwork and A.~Roth, ``The algorithmic foundations of differential privacy,'' \emph{Foundations and Trends in Theoretical Computer Science}, vol.~9, no.~3--4, pp.~211--407, 2014.

\bibitem{prasser2014arx}
F.~Prasser, F.~Kohlmayer, R.~Lautenschl\"{a}ger, and K.~A. Kuhn, ``{ARX}---A comprehensive tool for anonymizing biomedical data,'' \emph{AMIA Annual Symposium Proceedings}, pp.~984--993, 2014.

\bibitem{wilson2020dpsql}
R.~J. Wilson, C.~Y. Zhang, W.~Lam, D.~Desfontaines, D.~Simmons-Marengo, and B.~Gipson, ``Differentially private {SQL} with bounded user contribution,'' \emph{Proceedings on Privacy Enhancing Technologies}, vol.~2020, no.~2, pp.~230--250, 2020.

\bibitem{templ2017sdc}
M.~Templ, \emph{Statistical Disclosure Control for Microdata: Methods and Applications in {R}}.\hskip 1em plus 0.5em minus 0.4em Springer International Publishing, 2017.

\bibitem{apple2017dp}
{Apple Differential Privacy Team}, ``Learning with privacy at scale,'' \emph{Apple Machine Learning Journal}, Dec.\ 2017.

\bibitem{amer2025certified}
O.~Amer, S.~Chakrabarti, K.~Chakraborty, S.~Eloul, N.~Kumar, C.~Lim, M.~Liu, P.~Niroula, Y.~Satsangi, R.~Shaydulin, and M.~Pistoia, ``Applications of certified randomness,'' \emph{Nature Reviews Physics}, 2025.

\end{thebibliography}

\end{document}
